\begin{blocksection}
\question Identify the number of calls to \lstinline$eval$ and the number of calls
to \lstinline$apply$ in \lstinline{evaluator.gleam} made when \lstinline{run} calls
\lstinline{draw} in each program below. (\lstinline{run} evaluates the body of
\lstinline{draw} directly, so \lstinline{draw} itself is not an \lstinline$apply$.)

\begin{parts}
\part \begin{lstlisting}
fn draw(shape_slider: Int) -> Int {
  case True {
    True -> abs(3)
    False -> abs(1)
  }
}
\end{lstlisting}

\begin{tabular}{lrrrr}
\lstinline$eval$ & 4 & 5 & 7 & 9 \\
\lstinline$apply$ & 1 & 2 & 3 & 4
\end{tabular}

\begin{solution}
7 \lstinline$eval$, 1 \lstinline$apply$
\newline
Evals: (1) on the \lstinline{At} around the body, (2) on the entire \lstinline{case}, (3) on the subject \lstinline{True} (the \lstinline{True} before \lstinline{->} is a pattern, which is only compared against the subject's value, never evaluated; the other arm is never evaluated either), (4) on the \lstinline{At} around the arm's body, (5) on \lstinline{abs(3)}, (6-7) on \lstinline{abs}/\lstinline{3}
\newline
Apply: (1) \lstinline{abs} on \lstinline{3}
\end{solution}

\vspace{0.75\baselineskip}
\part \begin{lstlisting}
fn draw(shape_slider: Int) -> Point {
  case abs(2) {
    2 -> Point(1, 2)
    _ -> Point(0, 0)
  }
}
\end{lstlisting}

\begin{tabular}{lrrrr}
\lstinline$eval$ & 6 & 8 & 10 & 12 \\
\lstinline$apply$ & 1 & 2 & 3 & 4
\end{tabular}

\begin{solution}
10 \lstinline$eval$, 2 \lstinline$apply$
\newline
Evals: (1) on the \lstinline{At} around the body, (2) on the entire \lstinline{case}, (3) on \lstinline{abs(2)}, (4-5) on \lstinline{abs}/\lstinline{2}, (6) on the \lstinline{At} around the first arm's body (patterns are matched, not evaluated), (7) on \lstinline{Point(1, 2)}, (8-10) on \lstinline{Point}/\lstinline{1}/\lstinline{2}
\newline
Apply: (1) \lstinline{abs} on \lstinline{2}, (2) \lstinline{Point} on \lstinline{1} and \lstinline{2}
\end{solution}

\vspace{0.75\baselineskip}
\part \begin{lstlisting}
fn twin(x: Int) -> Point {
  Point(x, x)
}

fn draw(shape_slider: Int) -> Point {
  twin(abs(3))
}
\end{lstlisting}

\begin{tabular}{lrrrr}
\lstinline$eval$ & 6 & 9 & 11 & 14 \\
\lstinline$apply$ & 1 & 2 & 3 & 4
\end{tabular}

\begin{solution}
11 \lstinline$eval$, 3 \lstinline$apply$
\newline
Evals: (1) on the \lstinline{At} around the body of \lstinline{draw}, (2) on \lstinline{twin(abs(3))}, (3) on \lstinline{twin}, (4) on \lstinline{abs(3)}, (5-6) on \lstinline{abs}/\lstinline{3}, then in the body of \lstinline{twin}: (7) on its \lstinline{At}, (8) on \lstinline{Point(x, x)}, (9-11) on \lstinline{Point}/\lstinline{x}/\lstinline{x}
\newline
Apply: (1) \lstinline{abs}, (2) \lstinline{twin}, (3) \lstinline{Point}
\end{solution}

\vspace{0.75\baselineskip}
\part \begin{lstlisting}
fn pair(x: Int, y: Int) -> Point {
  Point(x, y)
}

fn draw(shape_slider: Int) -> Point {
  pair(abs(3), case True {
    True -> 2
    False -> 0
  })
}
\end{lstlisting}

\begin{solution}[1in]
15 \lstinline$eval$, 3 \lstinline$apply$
\newline
Evals: (1) on the \lstinline{At} around the body of \lstinline{draw}, (2) on \lstinline{pair(abs(3), case ...)}, (3) on \lstinline{pair}, (4) on \lstinline{abs(3)}, (5-6) on \lstinline{abs}/\lstinline{3}, (7) on the entire \lstinline{case}, (8) on the subject \lstinline{True} (the \lstinline{True} before each \lstinline{->} is a pattern, not an expression: it is only compared against the subject's value, never evaluated), (9) on the \lstinline{At} around the matching arm's body, (10) on \lstinline{2} (the \lstinline{False} arm is never evaluated), then in the body of \lstinline{pair}: (11) on its \lstinline{At}, (12) on \lstinline{Point(x, y)}, (13-15) on \lstinline{Point}/\lstinline{x}/\lstinline{y}
\newline
Apply: (1) \lstinline{abs} on \lstinline{3}, (2) \lstinline{pair} on \lstinline{3} and \lstinline{2}, (3) \lstinline{Point} on \lstinline{3} and \lstinline{2}
\end{solution}
\end{parts}
\end{blocksection}

\begin{blocksection}
\begin{guide}
\textbf{Teaching Tips}
\begin{itemize}
	\item This has historically been a tricky concept for students. \lstinline$apply$ may come off as easier to understand so relate it to just applying functions to already-evaluated arguments for students.
	\item Remind students of what gets \lstinline$eval$'ed: every \lstinline{Call} (the whole call, its function expression, and each operand), variables, literals, \lstinline{case} subjects and the chosen arm's body, and \lstinline{let} values and bodies.
	\item The parser also wraps the body of every function (and each \lstinline{let} value and each \lstinline{case} arm's body) in an \lstinline{At} marker, which costs one \lstinline$eval$. Decide ahead of time whether you want students to count these; the solutions above do.
	\item Operators are ordinary calls: \lstinline{1 + 2} is a \lstinline{Call} of the built-in \lstinline{+}, so it costs 4 \lstinline$eval$s (the call, \lstinline{+}, \lstinline{1}, \lstinline{2}) and 1 \lstinline$apply$. Record constructors like \lstinline{Point} are built-ins too.
	\item Every \lstinline{Call} evaluates all of its operands. Only \lstinline{case} skips work, by evaluating just the arm whose pattern matches (patterns are matched, not evaluated). \lstinline{&&} and \lstinline{||} are rewritten by the parser into a \lstinline{case}, which is how they short-circuit.
	\item For a \lstinline{Builtin}, \lstinline$apply$ just calls its \lstinline{run} function. A \lstinline{UserFn} \lstinline$eval$s its body in a new frame. Higher-order built-ins like \lstinline{list.map} call \lstinline$apply$ on the function they are given, once per element.
	\item Be very conscious about not accidentally evaluating the expressions yourself when you are counting; that's the interpreter's job!
	\item It might help to count the expressions by \lstinline$apply$ groups, i.e. count all of \lstinline$eval$ for one \lstinline$apply$able group and then move onto the next.
	\item Consider referencing Josh's helpful \href{https://drive.google.com/file/d/1YH6eaOAzPx8MC0fO60bIbq9oZipC2lEd/view}{walkthrough video}.
\end{itemize}
\end{guide}
\end{blocksection}